\section{一つの根号積分を三つの置換で比べる}\label{節:根号積分の三方法}

ここまでの道具を，同じ積分に使ってみよう。違いは答えそのものではなく，途中のどの式が簡単になるかにある。

\begin{example}{三角・双曲線・代数による三通りの計算}
$\int\sqrt{1+x^2}\dd x$を，三角置換，双曲線関数の置換，根号を含む代数的な置換によって求めよ。
\end{example}

\subsection{三角置換では正割の三乗になる}

$x=\tan\theta$，$-\frac\pi2<\theta<\frac\pi2$と置く。この範囲では$\sec\theta>0$なので
\[
\sqrt{1+x^2}=\sec\theta,\qquad
\dd x=\sec^2\theta\dd\theta,\qquad
\int\sqrt{1+x^2}\dd x=\int\sec^3\theta\dd\theta.
\]
根号はなくなったが，三乗の積分が残った。$J=\int\sec^3\theta\dd\theta$と書き，$\sec\theta$を微分する側，$\sec^2\theta$を積分する側にして部分積分する。
\[
\begin{aligned}
J&=\sec\theta\tan\theta
-\int\sec\theta\tan^2\theta\dd\theta\\
&=\sec\theta\tan\theta
-\int\sec\theta(\sec^2\theta-1)\dd\theta\\
&=\sec\theta\tan\theta-J+\int\sec\theta\dd\theta.
\end{aligned}
\]
任意定数を最後にまとめるものとして，この等式を$J$について解く。前に求めた正割の積分から
\[
J=\frac12\sec\theta\tan\theta
+\frac12\log(\sec\theta+\tan\theta)+C.
\]
この範囲では$\sec\theta>|\tan\theta|$だから，対数の中は正である。$x$へ戻して
\begin{equation}\label{式:一と二乗の根号積分}
\int\sqrt{1+x^2}\dd x
=\frac12\left\{x\sqrt{1+x^2}
+\log\bigl(x+\sqrt{1+x^2}\bigr)\right\}+C
\end{equation}
を得る。三角置換だけで計算が終わるのではなく，その後に部分積分を使った。

%
\subsection{双曲線置換では二乗の公式を使う}

$x=\sinh u$と置くと，実数$x$全体に対して$u$が一つに決まり，
\[
\sqrt{1+x^2}=\cosh u,\qquad
\dd x=\cosh u\dd u
\]
となる。今回は微分も同じ$\cosh u$になるので，積分は二乗で済む。式\eqref{式:双曲線の二倍公式}より
\[
\begin{aligned}
\int\sqrt{1+x^2}\dd x
&=\int\cosh^2u\dd u\\
&=\frac12\int\{1+\cosh(2u)\}\dd u\\
&=\frac u2+\frac14\sinh(2u)+C\\
&=\frac12\{u+\sinh u\cosh u\}+C.
\end{aligned}
\]
$u=\log(x+\sqrt{1+x^2})$と戻せば，式\eqref{式:一と二乗の根号積分}と一致する。

%
\subsection{根号と元の変数の和を置く}

第三の方法では，関数の名前を使わずに
\[
t=x+\sqrt{1+x^2}
\]
と置く。この和を選ぶ理由は，差との積が$1$になるためである。
\[
\bigl(\sqrt{1+x^2}+x\bigr)
\bigl(\sqrt{1+x^2}-x\bigr)=1.
\]
従って$\sqrt{1+x^2}-x=\frac1t$であり，和と差を取れば
\[
x=\frac12\left(t-\frac1t\right),\qquad
\sqrt{1+x^2}=\frac12\left(t+\frac1t\right),\qquad
\dd x=\frac12\left(1+\frac1{t^2}\right)\dd t.
\]
$t>0$であり，この$x$の式は$t>0$から実数全体への狭義単調増加な対応である。根号も微分も有理式へ変わったので，
\[
\begin{aligned}
\int\sqrt{1+x^2}\dd x
&=\frac14\int\left(t+\frac2t+\frac1{t^3}\right)\dd t\\
&=\frac{t^2}{8}+\frac12\log t-\frac1{8t^2}+C.
\end{aligned}
\]
ここで
\[
x\sqrt{1+x^2}=\frac14\left(t^2-\frac1{t^2}\right)
\]
だから，再び同じ答えになる。

実は第二，第三の置換は直接つながっている。$x=\sinh u$なら
\[
t=x+\sqrt{1+x^2}=\sinh u+\cosh u=\e^u.
\]
つまり第三の方法は，双曲線の媒介変数$u$の代わりに$\e^u$を変数として，式を有理式へ直したものである。三角置換は円の関係式$1+\tan^2\theta=\sec^2\theta$を使い，双曲線置換は$\cosh^2u-\sinh^2u=1$を使う。同じ平方根に対して，どの恒等式と微分の組が短い計算を作るかが違う。

最後に答えを微分して確認しておく。$s=\sqrt{1+x^2}$と書けば$s'=\frac xs$であり，
\[
\frac{\dd}{\dd x}\log(x+s)
=\cfrac{1+\cfrac xs}{x+s}=\frac1s.
\]
従って式\eqref{式:一と二乗の根号積分}の右辺の微分は
\[
\frac12\left(s+\frac{x^2}{s}+\frac1s\right)
=\frac12(s+s)=\sqrt{1+x^2}
\]
となる。三つの計算は同じ原始関数へ到達した。

%
